%*******************************************************
% Abstract in English
%*******************************************************
\pdfbookmark[0]{Abstract}{Abstract}


\begin{otherlanguage}{american}
	\chapter*{Abstract}
	
	Software forms the backbone of modern society, supporting critical infrastructure such as communication networks, financial systems, and healthcare services. The security of these systems is a pressing concern, as vulnerabilities can lead to severe consequences, including financial losses, reputational harm, and even threats to human safety. While software vendors are motivated by economic incentives and legal obligations—such as compliance with regulations like the EU's \ac{GDPR}—to prioritize security, the rapidly evolving threat landscape presents significant challenges. Emerging threats like advanced persistent attacks, ransomware, and supply chain vulnerabilities require more dynamic and scalable security solutions than traditional methods can offer.

	This thesis investigates the potential of \ac{AI}-driven penetration testing, specifically through the use of \ac{LLM}s, to automate and improve the process of identifying software vulnerabilities. The study begins by situating penetration testing within the \ac{SDLC}, highlighting its role as a creative and adaptive method for simulating real-world cyberattacks. It then explores the transition from manual testing to \ac{AI}-assisted approaches, analyzing recent developments such as PentestGPT, PentestAgent, and ARTEMIS. These advancements demonstrate that \ac{LLM}s, when combined with domain-specific knowledge, can effectively perform penetration tests in controlled settings.

	However, real-world applications reveal certain limitations. Tools like XBOW, despite their success in bug bounty programs, often generate high rates of false positives and struggle with proprietary or niche technology stacks. Security experts note that while \ac{AI}-driven tools excel at identifying known vulnerabilities, they lack the lateral thinking and contextual adaptability that human testers provide. Additionally, security risks such as data leakage through prompt injection and biases in training data pose further challenges.

	The thesis concludes that \ac{AI} can substantially accelerate penetration testing, but it is not yet a substitution for human expertise. Future research should focus on enhancing \ac{AI}'s ability to navigate proprietary systems, reducing false positives, and addressing security risks. By overcoming these challenges, \ac{AI}-driven penetration testing can become a dependable and scalable solution, complementing and ultimately strengthening security efforts.
\end{otherlanguage}
